%% ch-meta-limits.tex -- Part 2 (\ref{part:meta}), chapter 4: what is not claimed, and what is open.
\providecommand{\DS}[1]{\hyperref[ds:#1]{DS-#1}}

\chapter{Limits of the self-verification}
\label{ch:meta-limits}

%% =========================================================
\section{What is not claimed}\label{sec:meta-not-claimed}

\begin{itemize}
\item No absolute consistency statement of any kind (\DS{2}).
\item No formal soundness theorem for the core (K1).  Within VNB, K1 could
  be a theorem only relative to set models: VNB cannot define truth over
  the class universe, and cannot prove that a model exists.  The
  soundness of the core for version-X is link~3 of
  Section~\ref{sec:meta-chain}, a hand-written argument (\DS{6}).
\item No formal statement that a checker computes its relation (K3).  The
  transcription is by hand (\DS{9}).
\item Nothing about the seven oracle relations, \texttt{ineq},
  \texttt{arith-ground}, \texttt{arith-forsome}, \texttt{arith-simplify},
  \texttt{comm-ring-simplify}, \texttt{ring-simplify} and \texttt{sos}
  (\DS{5}).  They are checked -- \texttt{ineq} and \texttt{sos} by a
  certificate the checker verifies by arithmetic, the others by an
  independent second decision (Section~\ref{sec:oracles}) -- but the
  checks add no proving power and are not a proof of the oracle.
\item Nothing about the certificate store and its hash, \texttt{qed} and
  the install, the undo journal, the Scheme of the expression layer
  (\texttt{alpha-equiv?}, \texttt{subst-free}, the matchers), or MIT
  Scheme.  The encoded operations restate the expression layer; they do
  not verify it.
\item The META theorems are VNB theorems with empty bills, but most of
  their certificates list the oracles \texttt{arith} or \texttt{ineq}.
  They rest on the oracle assumption as every theorem does.
\end{itemize}

%% =========================================================
\section{What is open}\label{sec:meta-open}

\subsection{Coverage of K2}

\begin{itemize}
\item The open rows of lemma~(*): ground arithmetic and arithmetic on
  certified numbers, $\mathrm{LENGTH}$ of a term typed in a class of
  tuples through a \texttt{sep} domain or a functoid body, an applied
  structure operation, and the view clause
  (Section~\ref{sec:meta-lemma-star}).  The B* rows are proven for every
  certificate built from the other rows; a certificate that uses an open
  clause is not covered.
\item The rewriting rules (PS) are proven at any finite set of positions
  (Section~\ref{sec:meta-positions}).  The \texttt{lambda-beta} licence
  through its binder's own scope, the binder-domain reads of the list
  definedness checks, a macete with a chain of conditions, and the
  numeral and successor bridge of the macete matcher are not covered.
\item The certified omissions of the rules that owe the definedness of a
  class (Section~\ref{sec:meta-owed}): the figure in which the kernel
  omits the owed sequent because its certificate accepts the class is
  open for \texttt{union-intro} at any entry and for
  \texttt{lambda-type} at a binder list.  For the empty tuple and for
  \texttt{lambda-type} at one binder it is covered when one member of
  the context certifies the class, not for a composition of the rows.
\end{itemize}

\subsection{Faithfulness (K3)}

K3 is deferred until K2 is complete (\DS{9}).  Its obligations:
\begin{itemize}
\item that each actual kernel step is a member of its transcribed
  relation, and that each C row is faithful (named syntax against locally
  nameless; eigenvariable and capture conditions);
\item the rows of $\mathrm{FUN\text{-}DEF}$, with the lambda-lifting of
  records, and of the predicate table $\mathrm{PRED\text{-}DEF}$, which has
  no row yet; until it has, every definition in $T$ of a predicate without
  a row must have a predicative body;
\item the head codes still missing (\texttt{pair}, \texttt{singleton},
  \texttt{power}, \texttt{choice}, the arithmetic heads, the structure
  operations);
\item that $T$ contains the encoded catalog axioms the induction and
  class-extensionality theorems assume;
\item that when the kernel certifies a term through clauses that are
  rows of lemma~(*), its definedness sequent is among the sequents the
  rows reach (Section~\ref{sec:meta-lemma-star});
\item that the kernel's predicativity test of a \texttt{comp} condition
  decides exactly $\mathrm{NBG\text{-}OK}$, the premise of def-comp;
\item the encoding of a destructuring lambda over a product, with the
  two properties argued for it: every member of a product is the list of
  its projections, and the encoded body is predicative exactly when the
  body is (Section~\ref{sec:meta-owed});
\item that the chain of one-position rewrites, in the order of
  Section~\ref{sec:meta-positions}, is the relation the rewriting
  checkers enforce;
\item that a value the kernel's value test accepts as a set has its
  sethood derivable, as $\mathrm{J\text{-}LB}$ requires;
\item the independence from the builder's certificate of the checkers
  that re-run it (Section~\ref{sec:meta-base}; \DS{32}).
\end{itemize}

\subsection{Hand-written arguments}

\begin{itemize}
\item The soundness of each core rule for version-X, and the truth of
  each axiom of the catalog, under the readings of \DS{4} and \DS{7}
  (link~3), are argued in Chapter~\ref{ch:meta-soundness} and, for
  functoids, in Chapter~\ref{ch:meta-functoids}.  The argument is
  conditional on the hypotheses (S1)--(S5) of
  Section~\ref{sec:meta-sound-hyp}.
\item The definitions of the library are covered by their general form
  (definitional equations of new heads, guarded definitions read as
  descriptions), not one by one.
\end{itemize}

\subsection{Enforcement}

\begin{itemize}
\item The macete builder offers a rewrite at a record variable, which the
  checker refuses.  The disagreement is in the safe direction.
\end{itemize}
